\chapter{Design Methodology}\thispagestyle{empty}
\graphicspath{{Chapter3/figures/}}
{\setstretch{1.5}

\section{System Design Overview}

The LungInsight system design follows a comprehensive methodology integrating data engineering, machine learning architecture, and clinical requirements. The design process involved iterative consultation with domain experts, analysis of existing CAD systems, and evaluation of available datasets and technological options. The methodology emphasizes modularity, scalability, interpretability, and clinical usability.

\section{System Architecture}

The LungInsight architecture comprises six integrated layers:

\begin{itemize}
    \item \textbf{Data Layer:} Manages medical image data, metadata, and auxiliary clinical information. Implements secure storage with DICOM support and access controls.
    \item \textbf{Preprocessing Layer:} Handles image normalization, quality assessment, and artifact removal. Implements standardization protocols ensuring consistent input representation.
    \item \textbf{Feature Learning Layer:} Contains pre-trained and fine-tuned deep learning models for automatic feature extraction from images.
    \item \textbf{Classification Layer:} Applies ensemble methods combining multiple model predictions for robust disease classification.
    \item \textbf{Interpretability Layer:} Generates visual explanations (attention maps, Grad-CAM) highlighting diagnostic regions influencing model decisions.
    \item \textbf{Application Layer:} Provides user interfaces, API endpoints, result management, and integration with healthcare IT systems.
\end{itemize}

\section{Requirements Analysis}

\subsection{Functional Requirements}

\begin{enumerate}
    \item Accept medical images in multiple formats (DICOM, PNG, JPG, TIFF).
    \item Perform automatic image quality assessment and standardization.
    \item Detect and classify tuberculosis, bacterial pneumonia, viral pneumonia, COVID-19 pneumonia, and normal conditions.
    \item Generate classification probability scores for each disease category.
    \item Produce visual attention maps highlighting diagnostic regions.
    \item Store results with timestamps, image metadata, and audit trails.
    \item Provide API endpoints for integration with existing healthcare systems.
\end{enumerate}

\subsection{Non-Functional Requirements}

\begin{enumerate}
    \item \textbf{Performance:} Process chest X-rays in under 5 seconds per image on standard server hardware.
    \item \textbf{Accuracy:} Achieve sensitivity \textgreater 90\% and specificity \textgreater 85\% for disease detection.
    \item \textbf{Scalability:} Support concurrent processing of multiple images; horizontally scalable architecture.
    \item \textbf{Security:} Implement HIPAA-compliant data handling, encrypted storage, and access controls.
    \item \textbf{Reliability:} System uptime target of 99.5\% with graceful degradation.
    \item \textbf{Interpretability:} All predictions include visual explanations understandable by radiologists.
\end{enumerate}

\section{Technology Selection}

\subsection{Programming Language and Frameworks}

\textbf{Python} was selected as the primary language due to: (1) extensive deep learning ecosystem (TensorFlow, PyTorch, Keras); (2) rich data science libraries (NumPy, Pandas, scikit-learn); (3) wide adoption in medical AI research; (4) ease of prototyping and deployment.

\subsection{Deep Learning Framework}

\textbf{TensorFlow with Keras API} was chosen for: (1) production-ready deployment capabilities; (2) extensive pre-trained model zoo; (3) superior performance optimization tools; (4) strong community support; (5) better integration with cloud platforms.

\subsection{Image Processing Libraries}

\textbf{OpenCV} and \textbf{scikit-image} provide comprehensive image manipulation capabilities: image resizing, normalization, histogram equalization, filtering, and morphological operations.

\subsection{Web Framework}

\textbf{Flask or FastAPI} serves as the application framework for: (1) lightweight, flexible REST API development; (2) easy integration with machine learning pipelines; (3) minimal overhead for inference serving.

\subsection{Containerization}

\textbf{Docker} containerization ensures: (1) consistent deployment environments; (2) easy scaling and orchestration; (3) isolation from system dependencies; (4) reproducibility across development and production.

\section{Model Architecture Design}

\subsection{Architecture Selection}

The design employs multiple CNN architectures in an ensemble framework:

\begin{itemize}
    \item \textbf{DenseNet-121:} Compact architecture (7.98M parameters) balancing performance and computational efficiency. Dense connections enable efficient feature reuse.
    \item \textbf{ResNet-50:} Well-established architecture (25.5M parameters) with residual connections enabling training of very deep networks.
    \item \textbf{EfficientNet-B2:} Optimized architecture (9.1M parameters) achieving state-of-the-art accuracy-efficiency trade-off.
\end{itemize}

Each model is initialized with ImageNet pre-trained weights and fine-tuned on medical imaging datasets.

\subsection{Ensemble Strategy}

Individual model predictions are combined using weighted averaging:

\begin{equation}
P_{ensemble} = w_1 \cdot P_{DenseNet} + w_2 \cdot P_{ResNet} + w_3 \cdot P_{EfficientNet}
\end{equation}

where $w_i$ are learned weights optimized through validation set performance, with $\sum w_i = 1$.

\section{Data Strategy}

\subsection{Dataset Selection}

The system is trained and validated on publicly available datasets:

\begin{itemize}
    \item \textbf{ChexPert:} Large-scale chest X-ray dataset with 224,316 images and labels for 14 thoracic pathologies.
    \item \textbf{COVIDx CXR-4:} Curated COVID-19 chest X-ray dataset with COVID-positive, COVID-negative, and other pneumonia images.
    \item \textbf{TB X-Ray Database:} Tuberculosis detection dataset with anterior-posterior and lateral radiographs.
    \item \textbf{RSNA Pneumonia Detection Challenge:} Dataset specifically curated for pneumonia detection with expert annotations.
\end{itemize}

\subsection{Data Preprocessing Pipeline}

\begin{enumerate}
    \item \textbf{DICOM Processing:} Extract images from DICOM files, convert to 8-bit grayscale, standardize orientation.
    \item \textbf{Normalization:} Apply min-max normalization to [0, 1] range and standardization using ImageNet statistics for consistency with pre-trained models.
    \item \textbf{Resizing:} Uniformly resize all images to 224×224 pixels to match model input specifications.
    \item \textbf{Quality Control:} Automatically flag images with extreme histogram properties, excessive noise, or artifacts.
\end{enumerate}

\subsection{Data Augmentation}

Training augmentation includes:

\begin{itemize}
    \item Random rotation (±10 degrees)
    \item Random horizontal flip (probability 0.5)
    \item Elastic deformations (to simulate natural anatomical variations)
    \item Intensity perturbations (brightness, contrast adjustments ±10\%)
    \item Mixup (weighted averaging of random image pairs)
\end{itemize}

\section{Model Training Strategy}

\subsection{Optimization Algorithm}

Adam optimizer with initial learning rate 0.001, momentum 0.9, and learning rate decay (step decay reducing by factor of 0.1 every 10 epochs).

\subsection{Loss Function}

Weighted cross-entropy loss addressing class imbalance:

\begin{equation}
L = -\sum_i w_i \cdot y_i \cdot \log(\hat{y}_i)
\end{equation}

where $w_i$ are class weights inversely proportional to class frequencies.

\subsection{Regularization}

\begin{itemize}
    \item L2 regularization (weight decay) with coefficient 1e-4
    \item Dropout with rate 0.3 on fully connected layers
    \item Batch normalization on all convolutional layers
    \item Early stopping monitoring validation loss with patience of 15 epochs
\end{itemize}

\subsection{Training Protocol}

\begin{itemize}
    \item Batch size: 32 (balanced between memory requirements and convergence properties)
    \item Maximum epochs: 100 (with early stopping typically terminating around epoch 50-70)
    \item Train-validation split: 80-20 stratified split preserving class distributions
    \item Cross-validation: 5-fold cross-validation for robust performance estimation
\end{itemize}

\section{Interpretability Design}

\subsection{Attention Map Generation}

Grad-CAM (Gradient-weighted Class Activation Mapping) is implemented to visualize which image regions contribute most to disease predictions:

\begin{equation}
\text{Attention}_{c}(i,j) = \sum_k \frac{\partial y_c}{\partial A_{k(i,j)}} \cdot A_k(i,j)
\end{equation}

where $y_c$ is the output for class $c$, $A_k$ is the $k$-th activation map, and the gradient term indicates contribution of each spatial location.

\subsection{Confidence Calibration}

Predicted probabilities are calibrated using temperature scaling to provide well-calibrated confidence scores reflecting true prediction certainty.

\section{System Integration Design}

\subsection{API Design}

RESTful API endpoints:

\begin{itemize}
    \item \texttt{POST /api/predict:} Submit image for analysis
    \item \texttt{GET /api/results/{id}:} Retrieve results for a specific analysis
    \item \texttt{GET /api/models:} List available model versions
    \item \texttt{POST /api/feedback:} Collect radiologist feedback on predictions
\end{itemize}

\subsection{User Interface}

Web-based interface featuring:

\begin{itemize}
    \item Image upload functionality with drag-and-drop support
    \item Visualization of input image and preprocessed version
    \item Display of prediction results with confidence scores
    \item Interactive visualization of attention maps
    \item Result history and comparison capabilities
    \item User management and access controls
\end{itemize}

\section{Validation and Testing Strategy}

\subsection{Internal Validation}

\begin{enumerate}
    \item Unit testing for individual components (image processing, model inference)
    \item Integration testing for API endpoints and inter-component communication
    \item Performance testing under various load conditions
    \item Security testing (SQL injection, file upload vulnerabilities, unauthorized access)
\end{enumerate}

\subsection{Clinical Validation}

\begin{enumerate}
    \item Comparison against radiologist annotations on held-out test sets
    \item Inter-observer agreement assessment between radiologists and models
    \item Sensitivity and specificity analysis for each disease category
    \item Analysis of failure cases and model limitations
\end{enumerate}

\subsection{Robustness Testing}

\begin{enumerate}
    \item Testing on images from different equipment manufacturers and protocols
    \item Evaluation on images with various quality levels and artifacts
    \item Assessment of performance across different anatomical presentations
    \item Stress testing with extreme case presentations
\end{enumerate}

\section{Development Methodology}

The project follows Agile principles with:

\begin{itemize}
    \item Sprint-based development (2-week sprints)
    \item Iterative prototyping and testing
    \item Regular review and refinement based on validation results
    \item Continuous integration and deployment practices
    \item Version control using Git with feature branches
\end{itemize}

\section{Quality Assurance}

Quality metrics include:

\begin{itemize}
    \item Code quality: Python style adherence (PEP 8), static analysis with pylint/flake8
    \item Test coverage: Minimum 80\% code coverage with unit and integration tests
    \item Performance benchmarks: Inference time \textless 5 seconds per image
    \item Accuracy metrics: Sensitivity \textgreater 90\%, specificity \textgreater 85\%
    \item Documentation: Comprehensive API documentation, usage guides, and technical documentation
\end{itemize}

}
